\subsection{Continuous actions and points of the completed universal cover}
\paragraph{Statement D0a (Proved; Part D(a)).}\label{res:D0a}
Let $U/K$ be a hyperbolic curve in characteristic zero, with a geometric base point. The inverse system of all connected pointed finite etale geometric covers defines $\widetilde U$, and their smooth proper normalizations define $X'$. The semilinear automorphisms of $\widetilde U$ that cover the given action on $U_{\overline K}$ form $E_U$. Give this group the topology of its actions on all finite-cover fibres, together with its action on each finite field of coefficients. With this topology, continuous sections $\Gamma_K\to E_U$ are exactly continuous semilinear actions of $\Gamma_K$ on this pointed-cover system compatible with the projection to $U$ and inducing the prescribed action on $\overline K$.

A fixed geometric pro-point at finite distance is a lift of a point of $U(K)$ and its stabilizer is the point section. A fixed pro-point at infinity is a lifted cusp whose inertia subgroup is normalized by the section. These statements give the dictionary requested in cote 149, section 3, manuscript pp.~29--30, scans 59--61.

\begin{proof}
For clarity, a geometric pro-point means a compatible family of geometric points of the finite levels. The fibre over a geometric point of $U$ is a simply transitive profinite $\pi$-set; the fibre over a boundary point in $X'$ is a profinite $\pi$-set with stabilizers $I_c$. Finite etale covers of $U$ with a fixed geometric fibre form its Galois category. An automorphism of the fibre functor is determined by its compatible permutations of the finite fibres. Requiring semilinearity and the given action on $U$ identifies the coefficient automorphism with its image in $\Gamma_K$. This constructs the extension $1\to\pi\to E_U\to\Gamma_K\to1$, including its profinite topology, and gives both directions of the asserted correspondence. It is not necessary to choose a presentation of $\pi$ for this construction.

For the equivalent rational-function description, each function and each cover is defined at a finite level over a finite coefficient field. An open stabilizer of the finitely many coefficients, functions and their conjugates determining that level is precisely a basic neighbourhood in the topology just described. Conversely, fixing a sufficiently refined finite pointed cover fixes every specified finite collection of its rational functions. Thus on this group of automorphisms over $U$ the inverse-limit topology is the topology of pointwise convergence on discrete rational functions and coefficients. Automorphisms without the condition of preserving the map to $U$ are addressed in Part D(e).

A fixed pro-point projects to a $\Gamma_K$-fixed geometric point of $U$, hence a $K$-rational point. Its arithmetic stabilizer maps isomorphically to $\Gamma_K$, because the geometric deck action is free there. Inclusion of the image of a section into that stabilizer is therefore equality. Conversely the decomposition section of a $K$-point fixes the corresponding lift. At a boundary pro-point, the arithmetic stabilizer is $D_c=N_{E_U}(I_c)$ by A2, so fixing it is exactly normalization of its inertia. All assertions are compatible with changing the geometric base point by an inner automorphism.
\end{proof}

\paragraph{Statement D0b (Proved; Part D(b), stronger uniqueness).}\label{res:D0b}
Let $F$ be a finite extension of a number-field cyclotomic closure, or a finitely generated extension of such a field. For every hyperbolic $U/F$ and every continuous section $s_F:\Gamma_F\to E_U$, the action on $X'$ has at most one fixed geometric pro-point in all. Consequently, for a number field $K$ and a section $s:\Gamma_K\to E_U$,
\[
 (X')^{s(\Gamma_K^\circ)}=(X')^{s(\Gamma_K)},
 \qquad \#(X')^{s(\Gamma_K^\circ)}\le1.
 \tag{D.1}
\]
The same assertion holds for finitely generated $K/\mathbb Q$. This proves all four requested uniqueness statements and the proper-curve assertion. It identifies the existence part of Conjecture B with the point-or-cusp surjectivity assertion, through D0a; no such surjectivity is assumed. The manuscript passages are scans 59--61, manuscript pp.~29--30.

\begin{proof}
Two distinct pro-points differ at a finite geometric covering level. Replace its subgroup by an open characteristic subgroup $H$ contained in it, and then refine $H$ so that the compactification of its covering has positive genus. Here are details of the latter step. A hyperbolic genus-zero curve has at least three punctures. Over the algebraic closure, carrying three of them to $0,1,\infty$ and pulling back the Fermat cubic gives an etale cover whose compactification has genus one. For positive genus no initial change is needed. A Galois closure followed by an open characteristic refinement gives a cover dominating this one; Riemann--Hurwitz preserves positive genus. Characteristic refinements exist since a finitely generated profinite group has finitely many subgroups of each bounded index. This is also the content of the opened Esnault--Hai Lemma 5.6, stated for a characteristic-zero genus-zero hyperbolic curve \cite{EsnaultHai2008}.

The group $H s_F(\Gamma_F)$ defines a finite etale $V/F$ with the two chosen pro-points projecting to distinct $F$-points $v_1,v_2$ of its compactification $Y$. They may be interior or boundary points. Both induce the same section of $\pi_1(Y)$: their stabilizers contain $s_F(\Gamma_F)$, whose image is a section after the compactification homomorphism. With $v_2$ as origin for $Y\to J_Y$, the difference class $[v_1-v_2]$ therefore has zero Kummer class. The finite multiplication covers on $J_Y$ identify this with divisibility by every integer in $J_Y(F)$. A1 and AE1 make it zero. A divisor of degree zero of the form $v_1-v_2$ can be principal for distinct points only on a genus-zero curve, since its function would have degree one. This contradicts the positive genus of $Y$ and proves uniqueness. The proof applies over every finite extension needed to define a level, and AE1 supplies the same arithmetic input over finitely generated extensions of $F$.

For (D.1), use the base-change identification $\pi_1(U_{K(\mu_\infty)})=E_U\times_{\Gamma_K}\Gamma_K^\circ$. The preceding argument applies to the restricted section. Normality of $\Gamma_K^\circ$ means that $s(\Gamma_K)$ preserves its fixed-point set. A singleton preserved by a group is fixed pointwise. Both sides of (D.1) are empty or the same singleton. This proves the equality also in the absence of a point, without making an existence assertion.
\end{proof}

In particular, two different interior points over a cyclotomic field cannot have the same full section class: choose compatible lifts for the common section and apply D0b. The exact abelianized obstruction and a finite cover separating the smallest example are supplied next. Over number fields, the smallest existence problem left by this argument is an arbitrary section for $\mathbb P^1_{\mathbb Q}\setminus\{0,1,\infty\}$; D0b settles uniqueness for it, while establishing neither an interior point nor a boundary point for that arbitrary section.
